% 图契约：本题的源、边界和推导逐项对应；箭头符号按图注解释。
\documentclass[tikz,border=7pt]{standalone}
\usepackage{fontspec,xeCJK,amsmath,amssymb}
\IfFontExistsTF{Times New Roman}{\setmainfont{Times New Roman}}{\setmainfont{TeX Gyre Termes}}
\IfFontExistsTF{Heiti SC}{\setCJKmainfont[AutoFakeBold=1.5]{Heiti SC}}{\setCJKmainfont[AutoFakeBold=1.5]{FandolHei-Regular}}
\usetikzlibrary{arrows.meta,bending,calc,patterns,positioning}
\definecolor{ink}{HTML}{243745}
\definecolor{blue}{HTML}{275DAB}
\definecolor{green}{HTML}{177C63}
\definecolor{red}{HTML}{BD4935}
\definecolor{orange}{HTML}{AA6C16}
\definecolor{purple}{HTML}{715393}
\tikzset{every node/.style={font=\fontsize{9.3}{12.5}\selectfont,text=ink},
axis/.style={-{Stealth[length=4pt]},line width=.65pt,draw=ink!65},
vec/.style={-{Stealth[length=5pt]},line width=1pt,draw=blue},
short/.style={-{Stealth[length=3pt]},line width=.8pt},
guide/.style={dashed,draw=ink!40,line width=.5pt},
card/.style={rounded corners=2pt,draw=blue!65,fill=blue!4,align=center,text width=6.5cm,inner xsep=9pt,inner ysep=7pt},
title/.style={font=\fontsize{10.5}{14}\selectfont\bfseries,text width=15cm,align=center},
label/.style={font=\fontsize{8.4}{11}\selectfont,align=center},
note/.style={label,text width=6.2cm}}
\begin{document}
\begin{tikzpicture}[x=1cm,y=1cm]
\path[use as bounding box] (-8,-3.5) rectangle(8,3.5);
\node[title] at (0,3.15) {M3-022a\quad 反射像源与透射等效源的适用区域};
\begin{scope}[shift={(-4,0)}]

\draw[ink,thick] (-2.3,0)--(2.3,0);
\fill[blue!5] (-2.3,-1.7) rectangle(2.3,0);
\draw[guide] (0,1.2)--(1.4,0)--(0,-1.2);
\fill[red](0,1.2)circle(2.5pt) node[right] {$Q$};
\draw[blue,dashed](0,-1.2)circle(3pt);\node[right] at(.1,-1.2){$Q^\prime$};
\fill[ink](1.4,0)circle(2pt) node[above right] {$A$};
\node[label] at(-1.5,1.6){a 区 $\varphi_a$};\node[label] at(-1.5,-1.5){b 区 $\varphi_b$};
\node[label] at(0,2.35){b 区的 $Q^{\prime\prime}$ 置于原电荷位置，另作等效};
\node[label] at(0,-2.35){在 $A$：$R_+=R_-$；分别匹配势与法向导数};
\end{scope}
\node[card,draw=blue!65,fill=blue!4] at (3.9,1.8) {$\varphi_a\propto Q/R_++Q^\prime/R_-$};
\node[card,draw=green!65,fill=green!4] at (3.9,0) {$\varphi_b\propto Q^{\prime\prime}/(\varepsilon_rR_+)$};
\node[card,draw=orange!65,fill=orange!4] at (3.9,-1.8) {两区的等效源分别使用\\不能把 $Q,Q^\prime,Q^{\prime\prime}$ 当作同一真实体系};
\end{tikzpicture}\end{document}